%!TEX root = ../template.tex
%%%%%%%%%%%%%%%%%%%%%%%%%%%%%%%%%%%%%%%%%%%%%%%%%%%%%%%%%%%%%%%%%%%%
%% chapter4.tex
%% NOVA thesis document file
%%
%% Chapter with lots of dummy text
%%%%%%%%%%%%%%%%%%%%%%%%%%%%%%%%%%%%%%%%%%%%%%%%%%%%%%%%%%%%%%%%%%%%

\typeout{NT FILE chapter4.tex}%

\chapter{Dittor - Anonymous and Sybil-Resilient Credential Scheme for Tor}
\label{cha:dittor}

O \textit{Dittor} é um esquema de credenciais descentralizado e resistente a \textit{sybils} que permite \textit{selective disclosure} de atributos de forma confidencial, anónima e \textit{unlinkable}, mesmo que uma das entidades emissoras de credenciais seja maliciosa, para ser incorporado no Tor. Inicialmente, surge da necessidade de prevenção de \textit{Sybil attacks} no Tor, motivado pelas técnicas demonstradas anteriormente e culminando num paradigma diferente para a mitigação deste tipo de ataque: ao invés de mecanismos para identificar possíveis \textit{sybils} na rede, é criado um sistema de credenciais que não permite que exista mais do que uma credencial para um dado pseudónimo de um utilizador. 

O \textit{Nym} projeta uma adaptação do \textit{Coconut} integrada com as \textit{Nym credentials} para que os seus utilizadores possam interagir com o sistema de forma verificada e anónima, enquanto o \textit{Dittor} surge como um esquema para acreditação dos \textit{relay operators}, de modo a que um utilizador voluntário que queira correr um nó na rede Tor terá que passar por este processo de validação de identidade para o poder fazer, garantindo uma ligação entre uma entidade e identidade única e exclusiva. A partir da credencial fornecida, este utilizador pode comunicar com as \textit{Directory Authorities} para assumir o controlo sobre um nó na rede.

Como o Tor preza acima de tudo a sua privacidade, surge a necessidade de criação de um sistema de autenticação descentralizado, de forma a que: a privacidade dos utilizadores não possa ser violada por uma entidade; não exista ligação possível entre a verificação e a entidade emissora, nem entre o requerente e o pedido; não exista ligação entre o utilizador e as suas credenciais, culminando num processo \textit{Zero-Knowledge} para todas as entidades participantes no sistema, incluindo as \textit{Directory Authorities}, e permitindo uma ligação entre os nós que uma entidade corre com uma identidade. O esquema de acreditação usado no \textit{Dittor} assemelha-se muito ao \textit{Coconut}, mais especificamente nestas propriedades, mas muda os mecanismos e acrescenta uma camada de proteção contra \textit{Sybil attacks}.

\section{Propriedades}
\label{sec:dittor_propriedades}

Antes de avançar para as mudanças técnicas, é necessário definir o objetivo do \textit{Dittor}. Para isto, associa-se um conjunto de propriedades que devem ser mantidas ao longo de toda a implementação deste projeto, que em muito se baseiam nas propriedades do \textit{Coconut} \cite{sonnino2018coconut}:

\begin{itemize}
    \item \textit{\textbf{Threshold authorities:}} É apenas necessário um subconjunto do conjunto de autoridades emissoras de certificados para criação de certificados parciais para o utilizador, onde para um número \textit{t} de autoridades, a emissão dos certificados tem complexidade temporal \textit{O(t)};
    \item \textit{\textbf{Blind issuance and Unlinkability:}} Existe uma modalidade de \textit{Zero-Knowledge} no processo de emissão das credenciais, em que as entidades responsáveis por fazê-lo não aprendem nenhuma informação relativa ao requerente. Para além disto, os certificados emitidos são \textit{unlinkable} entre eles, mesmo que todas as entidades participem num ataque;
    \item \textit{\textbf{Non-interactivity:}} As autoridades emissoras de certificados podem operar de forma independente entre elas, de forma a que não seja necessário nenhuma sincronização para a emissão de certificados;
    \item \textit{\textbf{Liveness:}} Esta propriedade relaciona-se com a primeira, visto que é suficiente que um \textit{threshold} de entidades emissoras de certificados se mantenham honestas para que o sistema continue em funcionamento;
    \item \textit{\textbf{Efficiency:}} Apesar de todas as propriedades de segurança, é necessário garantir que haja uma \textit{performance} adequada às expectativas dos utilizadores, de forma a que se torne num modelo útil para o dia-a-dia. Para que isto aconteça, é fulcral que as primitivas criptográficas e as provas de \textit{Zero-Knowledge} tenham um protocolo eficiente para todos os passos do \textit{Coconut}, de forma a que a complexidade temporal de \textit{O(1)} deste esquema se mantenha nesta adaptação;
    \item \textit{\textbf{Short Credentials:}} O número de elementos da credencial tem de ser independente do número de atributos nela contidos, para mitigar qualquer \textit{traffic correlation} que permita a associação entre credenciais, deanonimizando o requerente;
    \item \textit{\textbf{Sybil-Resilient:}} Esta propriedade é inovadora relativamente ao \textit{Coconut}, e visa a proteção contra \textit{Sybil attacks} da rede Tor, que é o objetivo deste projeto. Para que o leitor possa entender melhor o mecanismo de proteção, é necessário explicar este tópico com mais profundidade.
\end{itemize}

\section{Mitigação de Sybil Attacks}
\label{sec:dittor_sybil}

Apesar do esquema do \textit{Coconut} basear-se na validação de uma prova de \textit{Zero-Knowledge} sobre a satisfação de um predicado sobre um atributo para emissão de uma credencial \cite{sonnino2018coconut}, não existe nenhuma verificação sobre o requerente que previna que este assuma diferentes identidades para a mesma prova. Seria possível guardar um tipo de \textit{distributed ledger} com a informação sobre os requerentes e as \textit{Zero-Knowledge Proofs} fornecidas que os conectassem para impedir a reutilização das mesmas provas para diferentes identidades, de forma a ser conhecimento público entre as entidades emissoras de certificados, contudo isto revelaria problemas de segurança relativos à \textit{unlinkability} entre os certificados emitidos e os utilizadores que, mesmo que a informação fosse privada para as autoridades confiadas do Tor, poderia ser alvo de ataques e suscetível a fuga de dados. Para isto, surge uma mudança no esquema do \textit{Coconut} na validação das credenciais dos requerentes por parte das entidades emissoras, de forma a impedir que uma entidade possa assumir diferentes identidades.

\subsection{SyRA}
\label{sec:syra}

Em 2025, surgem as assinaturas \textit{SyRA (Sybil-resilient anonymous)}, que permitem assinaturas digitais assegurando \textit{threshold issuance}, \textit{unforgeability} (que impede que sejam forjadas assinaturas por utilizadores que não sejam o próprio), \textit{Sybil resilience}, preservação de privacidade e \textit{unlinkability} \cite{crites2025syra}. Estas assinaturas são usadas de forma a complementar uma \textit{verifiable credential} da identidade real do sujeito, que pode ser obtida a partir de qualquer identificador legal (como o Cartão de Cidadão português), biométrico (como a impressão digital) ou a partir de qualquer método legal capaz de identificação como o OAuth \cite{recordon2006openid}, por exemplo, com o objetivo de alta compatibilidade com qualquer entidade legal para emissão destas credenciais. Estas entidades responsabilizam-se por garantir que o certificado gerado pelas mesmas é válido e dentro dos moldes legais, garantindo que não existe mais do que uma identidade para o mesmo certificado, isto é, não existem dois sujeitos com o mesmo identificador civil, por exemplo. As assinaturas \textit{SyRA} são usadas com estas credenciais de forma a gerar um pseudónimos para cada contexto, ou seja para cada entidade legal emissora das credenciais, de forma a que não seja possível retirar informação que correlacione a identidade do utilizador com os pseudónimos, mesmo que todas sejam maliciosas. Estes pseudónimos são derivados da identidade única do utilizador, que quando associados com a chave pública da entidade responsável pela \textit{verifiable credential}, é criado um identificador deste utilizador para ser usado num sistema. É importante notar que o pseudónimo criado com cada entidade mantém-se igual independentemente do número de certificados emitidos, garantindo que o utilizador não se personifica por diferentes identidades. De seguida, observamos as diferentes fases para este modelo, que contemplam os passos para obter uma assinatura \textit{SyRA}.

A fase de \textbf{geração de chaves} engloba a verificação que cada conjunto de entidades emissoras de credenciais consegue inicializar a sua instância da assinatura \textit{SyRA}. Após receção de todas as confirmações de todas as entidades de cada conjunto, o requerente assina a assinatura \textit{SyRA} com uma chave de verificação, que fica guardada nesta.

A \textbf{emissão limitada de credenciais} inicia-se com a demonstração de que uma entidade é dona de uma prova de identidade \textit{s} de um requerente. Essa verificação é formalizada para um valor público \textit{xs}, que disponibiliza um acesso de informação partilhado entre a entidade dona da prova de identidade e cada entidade emissora de credenciais, culimando no processo de emissão das mesmas, que está limitado por um dado valor que define o número mínimo de autoridades necessárias para tal. Cada uma destas entidades possui a relação entre uma testemunha \textit{ws} e \textit{s}, que pode ser formalizada de forma a definir mais concretamente os pormenores de validação que devem ser executados entre a entidade e a parte verificadora, tendo sempre de haver a verificação da chave pública da primeira, que não têm fundamento em ser especificados de forma a que haja uma maior compatibilidade.

Existe um mapeamento que faz a ligação entre as entidades que comprovam a identidade de um utilizador com o pseudónimo gerado, de forma a que uma entidade maliciosa que tente comprovar uma identidade que já foi comprovada por uma entidade honesta ou que uma entidade honesta tente fazê-lo com um identificador que já foi associado a uma entidade maliciosa, fique com a ligação comprovada e guardada no mapeamento, prevenindo que se repita no futuro e tornando-se responsável pela mesma.

A \textbf{geração de assinaturas} pode ser requerida tanto por entidades honestas como maliciosas, para um contexto e uma mensagem, criando uma ligação com o adversário assim que é feito o pedido com um identificador associado e único para cada ligação. O adversário submete uma assinatura e um pseudónimo, juntamente com um conjunto de identificadores das entidades que verificaram a sua identidade que, caso possua múltiplas identidades e estas entidades sejam maliciosas, ficam caracterizadas como maliciosas no sistema, visto que cada entidade só pode ter um pseudónimo associado a cada utilizador. De seguida, o pseudónimo é verificado como único para a mesma chave de verificação do adversário, identidade e contexto.

A interface de \textbf{verificação} é usada para que entidades externas possam validar assinaturas e receber os seus pseudónimos associados. Para assinaturas geradas com as chaves de sessão das entidades verificadoras das credenciais, é analisada a caracterização das mesmas no sistema, para evitar entidades maliciosas, de resto o adversário analisa apenas o resultado da verificação. Para extrair o pseudónimo associado, a interface permite a submissão de assinaturas, que podem usar o mesmo para determinar na sua \textit{key-value store table} se já existe para um dado contexto. Caso a assinatura submetida seja inválida, é retornado um valor negativo como pseudónimo.

O conceito de \textbf{fuga de assinaturas} é mencionado \cite{crites2025syra} como trabalho futuro, origem da interface do adversário para receber as assinaturas de utilizadores que são corrompidos durante o protocolo, de forma a implementar primitivas criptográficas como \textit{verifiable random functions} \cite{david2024musen} para \textit{forward secrecy}.

\subsection{SASSI}
\label{sec:sassi}

No seguimento da função de assinaturas \textit{SyRA}, surge a construção de assinaturas anónimas e resilientes a \textit{Sybils}, \textit{SASSI (Sybil-resilient Anonymous Signatures with Stateless Issuing)}, que é provado executar o protocolo de assinaturas \textit{SyRA} de forma segura \cite{crites2025syra}. Existem duas partes principais do \textit{SASSI}:

\begin{enumerate}
    \item Emissão das chaves secretas de cada utilizador como uma função determinística da sua identidade única com a chave secreta conjunta das entidades emissoras;
    \item Geração de assinaturas anónimas e resilientes a Sybils por utilizadores a partir das chaves secretas obtidas.
\end{enumerate}

A chave secreta do utilizador é composta por duas provas \textit{VRF} de Dodis-Yampolskiy, que são \textit{VRFs} aplicadas a grupos bilineares, com tamanho constante para as provas e as chaves produzidas, que pode ser usada de forma distribuída e proativa \cite{dodis2005verifiable}, que se torna ideal para o seu efeito, tornando o protocolo de \textbf{emissão limitada de credenciais} uma vertente distribuída deste tipo de provas, garantindo que a entidade emissora não observa a identidade do utilizador e um ataque conjunto necessitaria de \textit{t} entidades emissoras, comprometendo a privacidade, daí ser limitada.

O \textit{SASSI} usa três primitivas criptográficas: distribuição secreta de segredos de Shamir, encriptação de ElGamal e as \textit{VRFs} de Dodis-Yampolskiy \cite{dodis2005verifiable}. A construção deste mecanismo tem quatro fases: a \textbf{geração de chaves}, a \textbf{emissão limitada de credenciais}, a \textbf{geração de assinaturas} e a \textbf{verificação}.

A fase de \textbf{geração de chaves} começa pela geração de chaves privadas secretas para partilha para cada entidade emissora de credenciais, um conjunto de chaves públicas e uma chave pública de verificação.

A \textbf{emissão limitada de credenciais} é um protocolo interativo entre uma entidade com um identificador real e t entidades emissoras, cada uma com a sua chave secreta, permitindo que estas possam gerar um par de chaves secretas correspondentes a duas provas \textit{VRF} de Dodis-Yampolskiy, sem produzir nenhum resultado para as mesmas. A partir disto, a entidade partilha segredos de forma verificável, a partir de \textit{broadcasting} para as entidades emissoras de credenciais, e gera uma \textit{ZNP} que prova que o partilhado por \textit{broadcast} é consistente com o identificador real, justificando a relação entre a entidade e a identidade. A entidade emissora das credenciais verifica então a \textit{ZNP}, que caso seja bem-sucedida usa valores aleatórios para criar um tuplo com a sua chave secreta e com o identificador do requerente, que são depois enviados para o mesmo. Após receber \textit{t} mensagens de \textit{t} emissores, o utilizador constrói as suas chaves secretas a partir da reconstrução das partilhas dos emissores, que depois usa para duas igualdades que caso sejam bem sucedidas, guarda as chaves e usa a credencial verificável, senão resulta em falha.

A \textbf{geração de assinaturas} é possível para uma entidade a partir do seu identificador, da chave secreta derivada do mesmo, do par de contexto e de mensagem e da chave de verificação do emissor, gerando uma assinatura num par de contexto e mensagem. A entidade verificadora verifica a assinatura a partir da construção do pseudónimo, sem revelar a sua chave secreta. Isto é possível a partir da encriptação da chave secreta com um elemento público, que de seguida deve provar que o resultado está bem formado relativamente à chave de verificação do emissor, usando os mesmos passos e provando que a construção está bem formada.

A fase da \textbf{verificação} começa para uma entidade verificadora quando recebe uma mensagem que inclui a chave de verificação do emissor de credenciais, um par de contexto e mensagem e uma assinatura, que é verificada e associada ao pseudónimo extraído.

\subsection{Peer DID}
\label{sec:peer_did}

O uso das assinaturas \textit{SyRA} para sistemas descentralizados é mencionado no final do projeto \cite{crites2025syra}, como trabalho futuro dentro da modalidade dos \textit{Peer DIDs}. A menção entra de acordo com a ideia de \textit{self-sovereign identity}, de forma a defender que os utilizadores em sistemas de DIDs devem poder controlar a sua identidade a partir deste tipo de assinaturas, usando o seu pseudónimo em benefício da sua autonomia digital. Para isto, é mencionado o uso de \textit{Peer DIDs}, que anula a necessidade de uma \textit{blockchain} pública para a partilha dos DIDs.

Os \textit{Peer DIDs} podem ser usados de forma a que para entrar em contacto com uma entidade pública ou utilizador, seja criado um \textit{pairwise DID}, que quando usado em conjunto com as assinaturas \textit{SyRA} para criação das credenciais que estão ligadas à sua identidade real, torna-se possível proteger a identidade do utilizador apesar de ser impossível relacionar as suas interações entre contextos em que surja, enquanto permite defesa contra \textit{Sybil attacks} .

A utilização das assinaturas \textit{SyRA} e o \textit{SASSI} motivaram este projeto, de forma a implementar um esquema de credenciais anónimas e resilientes a \textit{sybils} para o Tor.

\section{Detalhes da Implementação}
\label{sec:dittor_detalhes}

No seguimento da explicação sobre o funcionamento de \textit{ZKPs}, segue-se a implementação de \textit{non-interactive zero-knowledge proofs (NIZK)}, que contempla uma eficácia superior na fase de verificação relativamente a \textit{zk-SNARKs}, sem necessidade de um código de texto confiável comum, facilitando a verificação para utilizadores com baixa potência computacional que, neste contexto, é uma vantagem por motivar utilizadores a correr nós na rede Tor. Dentro do contexto da assinatura, o uso de provas matemáticas de conhecimento de logaritmos discretos e relação multiplicativa de expoentes permite simplificar a relação entre a emissão e a verificação da assinatura \cite{crites2025syra}.

A implementação da geração distribuída de chaves tem por base esquemas de assinaturas para emissão de credenciais anónimas de forma distribuída \cite{doerner2023threshold}, culminando num processo uniforme e imparcial de um canal de broadcast, uma rede síncrona e uma maioria honesta.